\section{Experimental Design}
\label{sec:design}

\subsection{Two factors, four arms}
\label{sec:arms}

Every arm sees the same tasks, the same warehouse, the same model, the same
answer-format instruction and the same scorer. The arms vary two factors
(Table~\ref{tab:arms}): whether the agent has the governed tools and the
instruction to consult them before writing SQL (\emph{scaffolding}), and
whether it is given the domain's semantics (\emph{content}).

\begin{table}[t]
\centering
\footnotesize
\setlength{\tabcolsep}{5pt}
\begin{tabular}{@{}lll@{}}
\toprule
 & no scaffolding & scaffolding \\
\midrule
no content & \armS{} & \armH{} \\
content    & \armM{} & \armC{} \\
\bottomrule
\end{tabular}
\caption{The four arms as a two-factor design (artifact names:
\texttt{schema\_only}, \texttt{contract\_hollow}, \texttt{manual\_prompt},
\texttt{contract}). \armH{} differs from \armC{} only in that its prose
fields are empty. \armM{} carries the same knowledge as \armC{} in its
original form, pasted into the prompt.}
\label{tab:arms}
\end{table}

\armS{} is the floor: three ungoverned tools (\texttt{list\_tables},
\texttt{describe\_table}, \texttt{execute\_sql}), a 281-character system
prompt, no documentation. It measures what the model can do from the schema
alone, and it doubles as our operational definition of model capability.

\armM{} is the published approach reimplemented. DABStep's own
\texttt{manual.md} is pasted verbatim into the system prompt, bringing it to
24{,}177 characters; the tools are identical to \armS{}. It validates the
harness against the leaderboard band, and it is the comparison a
practitioner actually faces: the alternative to a contract is not ignorance
but a long prompt.

\armC{} is the treatment: the nine governed tools over the frozen contract,
and a 2{,}475-character system prompt that names the domains and metrics
available and instructs the agent to look them up before writing SQL.

\armH{} is the control that makes the ablation an ablation: the same tools
and the same instruction over the hollow contract, with a 1{,}984-character
prompt. \armC{} differs from the baselines in tools, in a procedural
instruction and in content at once; \armH{} holds the first two fixed and
removes only the third.

Two qualifications on reading Table~\ref{tab:arms} as a factorial. The
content factor is the same knowledge in two forms, the manual itself and a
contract authored from the manual alone, so the \armM{}--\armC{} contrast
measures the form of delivery as well as its presence. And the scaffolding
factor changes the tool surface as well as the instruction, because that is
the bundle a curated catalogue actually brings; we measure the bundle, not
its parts.

\subsection{Four models}
\label{sec:models}

\begin{table}[t]
\centering
\footnotesize
\setlength{\tabcolsep}{3pt}
\begin{tabular}{@{}llllr@{}}
\toprule
\tabhead{Model} & \tabhead{Pinned id} & \tabhead{Route} & \tabhead{T} & \tabhead{Tasks} \\
\midrule
\mGLM{} & \texttt{glm-5.3-flash} & OpenRouter, fp8 & 0 & 401 \\
\mDS{} & \texttt{deepseek-v4-flash} & OpenRouter, fp8 & 0 & 279 \\
\mSON{} & Claude Sonnet~5 & Bedrock via gateway & --- & 401 \\
\mGPT{} & \texttt{gpt-5.6-sol} & OpenRouter & --- & 401 \\
\bottomrule
\end{tabular}
\caption{The four models, in the order used throughout: bare-schema hard
accuracy. Each is pinned to one provider endpoint with fallbacks disabled.
T is temperature; a dash means the route does not accept the parameter. \mDS{}'s
task count is its complete-case set (Section~\ref{sec:models}).}
\label{tab:models}
\end{table}

We ran the full design once on each of four models
(Table~\ref{tab:models}), each pinned to a single provider endpoint with
fallbacks disabled so that quantization, price and throughput are fixed
within a run. We call \mGLM{} and \mDS{} the two flash models and \mSON{}
and \mGPT{} the two frontier models; that is the vendors' tiering, not a
measured quantity. Models appear throughout in order of bare-schema hard
accuracy (13.9\%, 22.6\%, 22.9\%, 37.0\%), the only capability axis declared
before any result was read; Section~\ref{sec:threats} records where that
axis disagrees with the contract arm's ordering. The four runs are separate
experiments, not replicates: they differ in model, in provider and in one
library fix (\appref{app:harness}).

Reasoning effort is \texttt{medium} throughout, a nominal setting rather
than an equated one: under it the contract arm reasons for about 350 tokens
per task on \mGPT{} and about 2{,}050 on \mSON{}. Temperature is 0 on the
two flash models; \mGPT{} does not accept the parameter and \mSON{}'s
gateway rejects both \texttt{temperature} and \texttt{seed}, so neither
frontier run has a determinism control. \mSON{} is reached through an
enterprise gateway on Anthropic's Messages API with cache control, because
the gateway's OpenAI-compatible route bills every input token fresh at a
multiple that differs by arm, a caching confound.

Two departures from the pre-registered design are recorded here. The
pre-registered primary model,
\texttt{deepseek-v4-pro}, was never swept: a 50-task probe of the
bare-schema arm showed it does not differ from the flash tier (27.5\%
against 25.0\%, $p{=}1.0$), so it would not have supplied the capability
contrast it was chosen for; \mGPT{} was substituted by the same criterion,
\mSON{} was added later, and every contrast involving them is exploratory
(\appref{app:harness}).
And the \mDS{} run lost 29\% of its rows to provider rate limiting,
near-uniformly across arms and with 114 of the 122 lost tasks losing all
four arms together, so we report it on the 279 tasks scoreable in all four
arms (its complete-case set), with intervals about 20\% wider
(\appref{app:truncation}).

\subsection{What is held fixed}
\label{sec:fixed}

Within a run everything else is fixed across the four arms: row limit (50
rows per result), tool-call cap (40), token budget, temperature, reasoning
effort, model, endpoint, quantization, task set, gold set, scorer and the
verbatim answer-format instruction. When an agent exhausts its output
budget the harness forces a final answer rather than discarding the row;
the rows so rescued fell on the ungoverned arms (on \mGLM{}, 41 in \armS{},
18 in \armM{}, none in either governed arm), so the policy works against the
treatment. Three harness properties treated the arms unequally
(\appref{app:harness}): a validator defect on \mGLM{} that rejected
\texttt{COUNT(*)} as if it were \texttt{SELECT *} in the governed arms only,
so \mGLM{}'s contract figures are a floor on that axis; a 100-row preview
cap that reached the governed arms on 31 calls, in the treatment's favour;
and different provider routes across models, which makes cross-run cost
differences suggestive only.

\subsection{Tasks, golds and grading}
\label{sec:golds}

DABStep withholds golds for its 450 test tasks. We reconstruct them from the
leaderboard's published per-task files by plurality over answers the
benchmark's own grader marked correct, with a 75\% plurality threshold
(median realised share 0.981), which yields golds for 401 of 450 tasks (332
hard, 69 easy) after five manually falsified consensus answers are excluded.
The three largest question families, 59 tasks, were recomputed directly from
the database using the manual's matching rule and reproduce exactly
(\appref{app:golds}).

The reconstruction was then measured rather than argued. Both arms of the
main contrast, \armC{} and \armM{} on \mGLM{}, were swept over all 450 tasks
and submitted to the leaderboard, which graded them against the withheld
golds: 51.9\% against 18.8\% on the 378 hard tasks, paired McNemar
$p{=}4.9\times10^{-28}$. On this model the main contrast is therefore graded
by the benchmark, not by us. On the 401 shared tasks the two gradings agree
on 382 (95.3\%); all 19 disagreements are wrong values our golds concealed,
all lie on the hard split, and all run the same way, so our hard-split
figures run about 5.7 points high for \armC{} and 3.0 for \armM{}. The
excluded tasks score like the kept ones under official grading, and a proxy
calibrated on these submissions bounds the leniency in every other arm and
model at under 3 points on any contrast (\appref{app:leniency}).

Answers are graded by DABStep's own scorer, vendored verbatim at a pinned
revision; our own earlier normaliser was stricter than the benchmark, and on
a 12-task pilot it manufactured a significant result in the contract's
favour that the official rules do not support.
Every arm sees every task, so all comparisons are paired and use McNemar's
exact test~\cite{mcnemar,dietterich}, reported with the discordant-pair
count and its split; proportions carry Wilson
intervals~\cite{wilson,brown-cai-dasgupta}; no multiplicity correction is
applied, and Section~\ref{sec:threats} says which results that exposes.
Tables report scored accuracy, correct over correct plus incorrect with
harness failures excluded; the released results carry strict accuracy, which
counts every non-correct row as wrong, beside it.
